\section{Single Element in a Sorted Array}
\label{sec:bs-single-element}

\problemstatement{We are given a sorted array in which \textbf{every
element appears exactly twice}, except for one element that appears
\textbf{exactly once}. Find that single element, in $O(\log n)$ time.}

\begin{patternbox}
At first glance this looks like it has nothing to do with sortedness --
``find the element that appears once'' sounds like a job for hashing or
XOR, both $O(n)$. The keyword that should stop you is \emph{``sorted
array''} appearing in the same breath as a uniqueness condition: whenever a
problem about counts, pairs, or uniqueness also guarantees sortedness,
suspect that sortedness is providing exploitable \emph{positional}
structure -- here, specifically, that pairs of equal elements occupy
predictable (even, odd) index pairs \emph{before} the single element, and
that this pairing pattern shifts by one position immediately
\emph{after} it.
\end{patternbox}

\subsection*{Understanding the problem carefully}

Suppose the single (unpaired) element sits at index $s$. For every pair
of equal elements before $s$, since each pair occupies two consecutive
array slots and there is an even number of elements total before $s$ (every
element before $s$ is properly paired), the first element of each such pair
lands on an \textbf{even} index, and its duplicate lands on the very next
(\textbf{odd}) index. That is: for $i < s$, pairs look like
$(\textit{arr}[2k], \textit{arr}[2k+1])$ with $\textit{arr}[2k] =
\textit{arr}[2k+1]$. After $s$, the single element has shifted the parity
of every subsequent pair by one position, so pairs there look like
$(\textit{arr}[2k+1], \textit{arr}[2k+2])$ instead.

\begin{ideabox}[Candidate Space and Predicate]
The candidate space is indices $0, \dots, n-1$ (restricted to even indices,
by construction). The predicate is
$P(\textit{even index } i) = (\textit{arr}[i] = \textit{arr}[i+1])$ --
\texttt{true} for every even index strictly before the single element, and
\texttt{false} from the single element's (even-aligned) position onward.
We binary search for the first even index where $P$ becomes \texttt{false}
-- that index holds the single element.
\end{ideabox}

\subsection*{Why forcing \textit{mid} to be even is necessary and correct}

The predicate $P$ above is only meaningful when evaluated at \emph{even}
indices, because the pairing pattern is defined relative to even/odd
positions, not relative to $\textit{mid}$'s raw value. So at each step, if
the computed midpoint $\textit{mid}$ happens to be odd, we adjust it down
by one (to the preceding even index) before testing $P$. After this
adjustment, the same monotonicity argument as ordinary binary search
applies directly: $P$ is true for an initial contiguous run of even
indices and false afterward, with the crossover landing exactly on the
single element's position.

\subsection*{Algorithm}

\begin{enumerate}
  \item Initialize $\textit{lo} \gets 0$, $\textit{hi} \gets n-2$
        (we only need to compare each candidate to its right neighbor).
  \item While $\textit{lo} < \textit{hi}$:
  \begin{enumerate}
    \item $\textit{mid} \gets \textit{lo}+(\textit{hi}-\textit{lo})/2$.
    \item If \textit{mid} is odd: $\textit{mid} \gets \textit{mid}-1$
          (force to even).
    \item If $\textit{arr}[\textit{mid}] = \textit{arr}[\textit{mid}+1]$:
          the single element is to the right; $\textit{lo} \gets
          \textit{mid}+2$.
    \item Else: the single element is at \textit{mid} or to the left;
          $\textit{hi} \gets \textit{mid}$.
  \end{enumerate}
  \item Return $\textit{arr}[\textit{lo}]$.
\end{enumerate}

\begin{algorithm}[H]
\caption{Single Element in a Sorted Array}
\begin{algorithmic}[1]
\Procedure{SingleElement}{\textit{arr}}
  \State $\textit{lo} \gets 0$, $\textit{hi} \gets n-2$
  \While{$\textit{lo} < \textit{hi}$}
    \State $\textit{mid} \gets \textit{lo}+(\textit{hi}-\textit{lo})/2$
    \If{\textit{mid} is odd}
      \State $\textit{mid} \gets \textit{mid} - 1$
    \EndIf
    \If{$\textit{arr}[\textit{mid}] = \textit{arr}[\textit{mid}+1]$}
      \State $\textit{lo} \gets \textit{mid} + 2$
    \Else
      \State $\textit{hi} \gets \textit{mid}$
    \EndIf
  \EndWhile
  \State \Return $\textit{arr}[\textit{lo}]$
\EndProcedure
\end{algorithmic}
\end{algorithm}

\subsection*{Dry run}

\dryrun{Take $\textit{arr} = [1,1,2,3,3,4,4,8,8]$ ($n=9$).}

\begin{itemize}
  \item $\textit{lo}=0,\textit{hi}=7$: $\textit{mid}=3$ (odd) $\to 2$.
        $\textit{arr}[2]=2$, $\textit{arr}[3]=3$. Not equal. $\textit{hi}=2$.
  \item $\textit{lo}=0,\textit{hi}=2$: $\textit{mid}=1$ (odd) $\to 0$.
        $\textit{arr}[0]=1$, $\textit{arr}[1]=1$. Equal. $\textit{lo}=2$.
  \item $\textit{lo}=2,\textit{hi}=2$: loop ends ($\textit{lo}=\textit{hi}$).
\end{itemize}

The single element is $\textit{arr}[2] = 2$, which is correct.

\subsection*{C++ implementation}

\begin{lstlisting}
#include <bits/stdc++.h>
using namespace std;

int singleElement(vector<int>& arr) {
    int n = arr.size();
    int lo = 0, hi = n - 2;

    while (lo < hi) {
        int mid = lo + (hi - lo) / 2;
        if (mid % 2 == 1) mid--;

        if (arr[mid] == arr[mid + 1]) {
            lo = mid + 2;
        } else {
            hi = mid;
        }
    }
    return arr[lo];
}

int main() {
    vector<int> arr = {1, 1, 2, 3, 3, 4, 4, 8, 8};
    cout << "Single element: " << singleElement(arr) << endl;
    return 0;
}
\end{lstlisting}

\begin{complexitybox}
\textbf{Time Complexity.} The search window halves each iteration (the
parity adjustment only ever shifts \textit{mid} by one, not affecting the
overall halving rate), giving
\[
  O(\log n).
\]
\textbf{Space Complexity.} $O(1)$ auxiliary space.
\end{complexitybox}

\begin{takeawaybox}
Sortedness can encode positional structure beyond simple ordering --
here, sortedness plus a pairing guarantee creates an exploitable
\emph{parity} pattern. Whenever a sorted-array problem also involves
counts, pairs, or duplicates, look for whether index parity (even/odd
position) is doing some of the work that the array's values alone do in
simpler binary search problems.
\end{takeawaybox}
